\section{Per-Dimension Analysis}
\label{sec:analysis}

This section populates five comparison tables across the seven
configurations enumerated in \S\ref{sec:configurations}. The
configurations are referred to throughout by their table-row label
(C1 through C7) before they are formally introduced in
\S\ref{sec:configurations}; readers wanting the full descriptions
should jump ahead. Each table draws claims at one of two confidence
levels: the row labeled C1 (Stellar~+~BLS12-381~+~PLONK) is
generated from a first-party benchmark harness whose methodology is
specified in Appendix~\ref{app:bench} and whose source lives in this
paper's source tree under \texttt{benchmarks/c1/}; the C1 deployment
under measurement comprises five Soroban contracts whose verifying
keys all derive from a single SRS reused from the Ethereum
Foundation KZG ceremony~\cite{eth-kzg-ceremony}, and the harness
reports per-contract numbers. The aggregate C1 row in each comparison table
summarizes the five-contract distribution; per-contract sub-rows are
available in the cost table where the variation across contracts is
load-relevant. Rows C2--C6 cite published numbers attributed inline;
row C7 is a hypothetical configuration whose cells are model
estimates derived in \S\ref{sec:migration}, included precisely to
make the host-function constraint legible. All USD figures are
anchored to 2026-04 fee-market and spot-rate conditions; reviewers
wanting contemporary numbers should re-run the harness and re-cite.

\subsection{Verification cost}
\label{sec:analysis-cost}

Table~\ref{tab:cost} summarizes verification cost. Cells are placed
at TBD pending the §4 measurement-pass refresh; the structural claims
the table is designed to support are stable, and so we describe them
in prose:

\begin{itemize}
  \item \textbf{Native pairing precompiles dominate the cheap end.}
    Configurations whose verifier reduces to a small number of
    pairings called via a host function (C1, C2) are one to two
    orders of magnitude cheaper than configurations whose verifier
    runs the same primitives in pure bytecode. The Stellar
    BLS12-381 pairing host function~\cite{stellar-cap-0059} and the
    Ethereum BN254 pairing precompile (EIP-197) are the components
    that dominate this regime.
  \item \textbf{Universal updatable schemes carry verifier overhead.}
    PLONK-family verifiers require multiple $\mathbb{G}_1$ and
    $\mathbb{G}_2$ multi-scalar multiplications beyond the pairing
    check; on chains without an MSM precompile, this raises baseline
    verification cost by an additional $2$--$5\times$ over Groth16.
    On Stellar specifically (C1) the overhead is attenuated by
    CAP-0059's $\mathbb{G}_1$ scalar-multiplication and
    $\mathbb{G}_1$ addition host
    functions~\cite{stellar-cap-0059} --- the MSM is decomposed into
    native scalar-mul calls --- but is not eliminated; the C1 row
    sits in the cheap regime, just at a constant-factor disadvantage
    relative to a hypothetical Groth16-on-Stellar configuration.
  \item \textbf{Transparent-setup configurations on
    non-FRI-friendly chains are dominated.} C7 (the hypothetical
    Stellar~+~Plonky3~+~FRI) would require a pure-WASM FRI verifier
    on Stellar, with verification cost dominated by Merkle-path
    hashing and field arithmetic that the BLS12-381 host functions
    do not accelerate. The configuration is included in the table to
    document the host-function gap, not as a viable production
    target.
\end{itemize}

\begin{table*}[t]
  \centering
  \caption{Per-verification cost. C1 sub-rows are the median per-tier
    \texttt{verify\_membership} fee captured on Stellar testnet from
    the production deployment (release \texttt{v0.0.5}, captured
    2026-05-02; one stroop $=10^{-7}$ XLM); the C1 \emph{aggregate}
    row is the unweighted median across the five contracts'
    proof-bearing entrypoints. C7 is measured first-party from the
    \texttt{wt-pq-fri-feasibility} branch on the same testnet
    (captured 2026-05-02); proof-bearing rows are not yet measurable
    end-to-end (no PCS layer tying FRI to an AIR), so only deploy and
    admin ops are reported. C2--C6 cite published numbers; USD is
    anchored to a representative 2026-05 spot rate. Where deploy and
    \texttt{verify\_membership} differ materially within one
    configuration, both are reported; \texttt{deploy} is a
    one-time cost, \texttt{verify\_membership} the per-call cost.}
  \label{tab:cost}
  \begin{tabular}{llrrr}
    \toprule
    & Configuration / op & Native unit & USD$^\star$ & Source \\
    \midrule
    \multicolumn{5}{l}{\emph{C1 --- Stellar + BLS12-381 + PLONK (this work, release \texttt{v0.0.5}, 2026-05-02)}} \\
    \quad C1 deploy (median, 5 contracts)              & stroops   & 29{,}287    & \$0.0012  & this work \\
    \quad C1 verify\_membership (median across 5)      & stroops   & 79{,}114    & \$0.0032  & this work \\
    \quad C1 create\_group (median, tier 0)            & stroops   & 176{,}466   & \$0.0071  & this work \\
    \quad C1 update\_commitment (sep-anarchy tier 0)   & stroops   & 125{,}871   & \$0.0050  & this work \\
    \quad \quad C1.1 sep-anarchy verify\_m (tier 0)    & stroops   & 79{,}107    & \$0.0032  & this work \\
    \quad \quad C1.2 sep-democracy verify\_m (tier 0)  & stroops   & 79{,}148    & \$0.0032  & this work \\
    \quad \quad C1.3 sep-oligarchy verify\_m (revert)  & stroops   & 78{,}983    & \$0.0032  & this work \\
    \quad \quad C1.4 sep-oneonone create\_group        & stroops   & 144{,}317   & \$0.0058  & this work \\
    \quad \quad C1.5 sep-tyranny deploy only           & stroops   & 29{,}454    & \$0.0012  & this work \\
    \midrule
    C2  & Ethereum L1 + BN254 + Groth16 (verify)         & ${\sim}220$K gas$^\dagger$ & ${\sim}\$0.50$ & \cite{pertsev-2019-tornado,semaphore-2024} \\
    C3  & Aztec + UltraPlonk (L2 verify, batched)        & {amortised L2}             & ${<}\$0.01$    & \cite{aztec-2024} \\
    C4  & Mina + Pickles (protocol-level)                & {built-in}                 & {n/a}$^\ddagger$ & — \\
    C5  & Polygon zkEVM + Plonky2 + KZG/BN254            & ${\sim}300$K gas (L1 batch) & {amortised}    & — \\
    C6  & STARK over Goldilocks on StarkNet              & {steg/L1 batched}          & {amortised}    & \cite{ben-sasson-2018-stark} \\
    \midrule
    \multicolumn{5}{l}{\emph{C7 --- Stellar + FRI-style hash-based (this work, feasibility branch, 2026-05-02)}} \\
    \quad C7 deploy                                     & stroops   & 17{,}196{,}133 & \$0.69    & this work \\
    \quad C7 set\_restricted\_mode (admin)              & stroops   & 12{,}291      & \$0.00049  & this work \\
    \quad C7 create\_group / verify\_m / update\_comm.  & \multicolumn{2}{c}{\emph{not yet measurable}$^{\sharp}$} & this work \\
    \bottomrule
  \end{tabular}
  \\[0.5em]
  \footnotesize
  $^\star$ USD anchored to XLM/USD = \$0.40 and ETH/USD = \$2{,}300
  spot, 2026-05 representative; reviewers wanting current numbers
  should re-anchor to spot rates at submission time.\\
  $^\dagger$ C2 native-unit estimate is the canonical Tornado-shaped
  Groth16 verifier on Ethereum L1 (one BN254 pairing precompile call
  + one G1 MSM); the precise cell tracks the Tornado V1.4 verifier
  cost~\cite{pertsev-2019-tornado}. \\
  $^\ddagger$ Mina's verification is protocol-level (every block
  carries its own Pickles proof); per-application USD is not the
  meaningful unit. \\
  $^\sharp$ The C7 verifier deploys and accepts a bench-scope FRI
  low-degree-test proof, but the prover side does not yet carry a
  batched-PCS layer tying FRI commitments to an AIR; the
  proof-bearing entrypoints' on-chain cost is bounded above by the
  estimated $1$--$2$ orders of magnitude over C1 derived in
  \S\ref{sec:cfg-c7}, but is not yet first-party measured. The
  deploy-cost ratio of $\sim$590$\times$ over C1 is, by contrast,
  measured directly.
\end{table*}


\subsection{Proof size}
\label{sec:analysis-size}

Table~\ref{tab:proof-size} ranks configurations by proof size.
Groth16 is uniquely small at constant-size 192~B with BLS12-381
compressed encoding, or 128~B with BN254 (C2)~\cite{groth-2016-onpairing}.
PLONK-family proofs (C1, C3) are $\sim$0.5--1~KB; the C1 row sits
at $\sim$700~B for a vanilla KZG-PLONK proof on BLS12-381 with
compressed encoding (9~$\mathbb{G}_1$ + 7~$\mathbb{F}_r$ elements
in the Gabizon--Williamson--Ciobotaru
construction~\cite{gabizon-2019-plonk}). Recursive schemes (C4,
Mina~+~Pickles) carry the recursion-frame overhead and land in the
10--20~KB range. FRI-based schemes (C5--C7) are the largest at
50--200~KB depending on the security level and FRI rate parameter;
the Plonky2 and Plonky3 lines achieve smaller sizes than classical
STARKs by aggressive use of small fields and recursion within the
proof.

\begin{table}[t]
  \centering
  \caption{Proof size. Pairing-based Groth16 proofs are constant-size
    (192~B with BLS12-381, 128~B with BN254); PLONK with KZG sits at
    $\sim$700~B on BLS12-381; recursive schemes grow into the
    multi-KB range; FRI-based proofs are the largest. The C1 row
    reports a single value because PLONK proof size depends on the
    construction and curve, not on the verifying-key family --- all
    thirty C1 verifying keys produce proofs of the same wire size.
    The C7 row reports the bench-scope FRI proof size measured on
    the feasibility branch (low-degree test alone); a production
    PCS-tied FRI proof would land in C6's band.}
  \label{tab:proof-size}
  \begin{tabular}{lr}
    \toprule
    Configuration & Proof size \\
    \midrule
    C1  & ${\sim}700$ B (PLONK / KZG over BLS12-381, compressed)$^{\ddagger}$ \\
    C2  & 128 B (Groth16 / BN254 compressed) \\
    C3  & ${\sim}0.5\text{--}1$ KB (UltraPlonk) \\
    C4  & ${\sim}10\text{--}20$ KB (Pickles, recursion overhead) \\
    C5  & ${\sim}50\text{--}100$ KB (Plonky2 + outer KZG wrap) \\
    C6  & ${\sim}50\text{--}200$ KB (FRI) \\
    C7  & ${\sim}8$ KB (bench-scope, FRI low-degree test only)$^{\sharp}$ \\
    \bottomrule
  \end{tabular}
  \\[0.5em]
  \footnotesize
  $^{\ddagger}$ Vanilla KZG-PLONK with the
  Gabizon--Williamson--Ciobotaru proof
  shape~\cite{gabizon-2019-plonk}; production variants
  (TurboPlonk, Plonkup, jellyfish-PLONK) are similar order of
  magnitude with constant-factor differences from custom-gate and
  lookup overhead.\\
  $^{\sharp}$ Measured against the feasibility-branch prover with
  parameters $\log n=6$, $3$ FRI layers, $8$ queries, blow-up $2$;
  a production-shape PCS-tied FRI proof at the same security level
  would land in the $50$--$200$~KB band of C6.
\end{table}


The implication for registry design is direct: configurations whose
proofs travel in transactions on a chain with per-byte fees pay for
proof size linearly. C5--C7's 50--200~KB proofs are economical only
on chains that subsidize ZK-rollup-style submissions; on a general
L1, they are not.

\subsection{Prover wall-clock}
\label{sec:analysis-prover}

Table~\ref{tab:prover-time} reports prover wall-clock for the
canonical advancement relation at a medium ring size (64 members for
systems with a member notion; ``one signal'' for Tornado-shaped
systems). Single-threaded baseline on commodity x86\_64; multi-thread
scaling is reported in Appendix~\ref{app:bench} for C1 and is cited
inline for the others.

The headline ordering is:
\begin{enumerate}
  \item Plonky2/3 (C5, C7) win on prover time at the cost of proof
    size, by virtue of FRI's amenability to GPU- and SIMD-accelerated
    NTT and Merkle-tree workloads on small fields.
  \item Groth16 (C2) is competitive in the medium-circuit regime
    where the relation is well-tuned but degrades at large ring
    sizes; the per-proof cost is dominated by the multi-scalar
    multiplication on $\mathbb{G}_1$.
  \item PLONK-family schemes with KZG commitments (C1, C3) introduce
    additional polynomial-commitment overhead and are typically
    $1.5$--$3\times$ slower than Groth16 at equivalent constraint
    counts. C1's reuse of the EF KZG SRS does not change the
    prover-time profile relative to a same-size project-specific
    KZG ceremony --- the SRS-reuse property is a trust posture, not
    a performance one --- but it does eliminate the per-relation
    Phase-2 ceremony cost that a Groth16 deployment of the same
    five-policy structure would carry.
  \item STARK over Goldilocks (C6) is competitive with Plonky2 in raw
    prover throughput; the difference is a constant factor that
    depends on the specific commitment scheme and the rate parameter.
\end{enumerate}

\begin{table}[t]
  \centering
  \caption{Prover wall-clock for the canonical advancement relation
    at the medium ring size (64 members for systems with a member
    notion). Single-threaded baseline on commodity x86\_64. The C1
    row is the median across the five contracts of the C1
    deployment; per-contract numbers are in
    \texttt{benchmarks/c1/results.json}. Cells marked TBD pending
    the §4 measurement-pass refresh.}
  \label{tab:prover-time}
  \begin{tabular}{lr}
    \toprule
    Configuration & Prover wall-clock \\
    \midrule
    C1  & TBD ms (median across 5 contracts; harness in App.\ \ref{app:bench}) \\
    C2  & TBD ms (cited) \\
    C3  & TBD ms \\
    C4  & TBD ms \\
    C5  & TBD ms \\
    C6  & TBD ms \\
    C7$^\dagger$ & TBD ms (model) \\
    \bottomrule
  \end{tabular}
\end{table}


\subsection{Trust assumptions}
\label{sec:analysis-trust}

Table~\ref{tab:trust} summarizes trust posture along three axes:
trusted-setup ceremony, validator set composition, and operationally
realized censorship resistance.

\begin{table*}[t]
  \centering
  \caption{Trust assumptions. ``Setup'' names the ceremony posture;
    ``Validators'' summarizes the anchor chain's decentralization;
    ``Censorship'' is the operationally-relevant dimension for
    transactions reaching inclusion.}
  \label{tab:trust}
  \begin{tabular}{llll}
    \toprule
    & Setup ceremony           & Validator set                & Censorship resistance     \\
    \midrule
    C1 & Universal SRS (EF KZG, ${\sim}141$K)$^{\ddagger}$ & Stellar Validators ($\sim$30) & SCP quorum, low censorship      \\
    C2 & Per-circuit MPC          & Ethereum L1 ($\sim$10$^6$ stakers) & Highest in surveyed set       \\
    C3 & UltraPlonk universal     & Aztec sequencer (centralized in alpha) & Sequencer-mediated      \\
    C4 & Transparent (Pickles)    & Mina ($\sim$300 block producers) & Moderate                      \\
    C5 & KZG universal updatable  & L2 sequencer + L1 anchor   & L1-anchored fallback              \\
    C6 & Transparent (FRI)        & Anchor-dependent          & Anchor-dependent                  \\
    C7 & Transparent (FRI)        & Stellar Validators ($\sim$30) & Same as C1                       \\
    \bottomrule
  \end{tabular}
  \\[0.5em]
  \footnotesize
  $^{\ddagger}$ C1 reuses the Ethereum Foundation KZG Summoning
  Ceremony~\cite{eth-kzg-ceremony}; the trapdoor-honesty assumption
  is one-of-${\sim}141{,}000$ contributors, the same setup that
  secures Ethereum's EIP-4844 blob commitments.
\end{table*}


The trust-assumption ranking is partial-order, not total: Ethereum L1
(C2) wins on validator decentralization but loses to Mina (C4) on
trusted setup; Mina wins on setup but loses to Ethereum on validator
count. The deployment recommendations in
\S\ref{sec:recommendations} navigate this partial order by deployment
profile.

\subsection{Post-quantum posture}
\label{sec:analysis-pq}

Table~\ref{tab:pq} classifies each configuration by the §3.4
post-quantum class. The headline is uncomplicated: every
pairing-based configuration (C1--C5, including the seemingly modern
Plonky2~+~KZG~over~BN254 in C5) is broken under Shor at zero
post-quantum security level; only the FRI-based configurations
(C6, hypothetical C7) carry positive Grover-resistant security. The
C5 case warrants precision: the post-quantum soundness of a deployed
configuration is the minimum of its constituent commitment schemes'
post-quantum soundness, and the outer KZG wrap dominates the
minimum. The inner FRI prover is not irrelevant to C5 --- it remains
post-quantum-secure as a sub-component --- but it is not the
bottleneck; an adversary capable of forging KZG commitments produces
accepted L1 verifications regardless of the inner FRI's properties.
Modernity in the literature is uncorrelated with post-quantum
readiness because the post-quantum decision is orthogonal to
construction novelty.

\begin{table}[t]
  \centering
  \caption{Post-quantum posture. Class labels follow the §3.4
    enumeration: (P) pairing-based, (E) classical-EC non-pairing,
    (H) hash-based / FRI. The ``Grover-resistant level'' column
    states the PQ security against a generic Grover attack on the
    underlying primitive; a Grover-resistant level of zero means
    Shor breaks the underlying assumption regardless of hash
    parameters. The hash-family row references the underlying
    collision-resistance assumption, which is the same
    standard-model assumption that
    NIST PQC stack-based signatures (FIPS~205 SLH-DSA) rely on,
    even though SLH-DSA is a signature scheme rather than a SNARK.}
  \label{tab:pq}
  \begin{tabular}{lll}
    \toprule
    Configuration & Class & Grover-resistant level \\
    \midrule
    C1  & P  & 0 (KZG binding reduces to DL on $\mathbb{G}_1$) \\
    C2  & P  & 0 \\
    C3  & P  & 0 (KZG on BN254) \\
    C4  & E  & 0 (Pasta DL) \\
    C5  & P  & 0 (KZG on BN254 dominates the inner FRI) \\
    C6  & H  & ${\sim}100$ bits (depending on hash) \\
    C7  & H  & ${\sim}64$--$100$ bits (Poseidon2 over BN254 $\mathbb{F}_r$) \\
    \bottomrule
  \end{tabular}
\end{table}


The careful caveat about hash-based configurations is that the
Grover-resistant security level depends on the underlying hash, not
on the FRI construction itself. Poseidon and Poseidon2 in their
standard parameter sets target $\sim$128 classical bits, which under
generic Grover speedup gives $\sim$64 post-quantum bits before
parameter doubling. STARK-friendly hashes such as Rescue-Prime are
typically parameterized for $\sim$100 PQ-bits at the cost of
constraint count. The PQ posture column is therefore a property of
the (SNARK $\times$ hash) joint choice, not of the SNARK alone.

\subsection{Synthesis}
\label{sec:analysis-synthesis}

Across all five tables, the configurations group into three regimes:
\emph{cheap pairing-based with native host functions} (C1, C2),
\emph{moderate-cost universal-updatable on chains with sequencers}
(C3, C5), and \emph{post-quantum-ready transparent} (C6 and the
blocked C7). The cheap-and-PQ-ready cell is empty along the
operationally-relevant censorship-resistance and decentralization
axes: C6 anchored on StarkNet does combine FRI-friendly host
functions with hash-based PQ security at competitive cost on the L2
itself, but StarkNet's centralized prover and sequencer place the
configuration outside the §7.3 censorship-resistance profile and
outside the validator-set tier of C2. The cell that is actually
empty is therefore \emph{cheap, PQ-ready, and decentralized to a
tier comparable with C1 or C2}; the open problems of
\S\ref{sec:open} make the precise shape of the gap explicit.
